% Cover letter for the Journal of Sports Economics submission.
% Standalone one-page letter; addressed to the editor (not blinded).
\documentclass[12pt,USenglish]{article}
\usepackage[utf8]{inputenc}
\usepackage[T1]{fontenc}
\usepackage[margin=1in]{geometry}
\usepackage{mathptmx}            % Times, matching the manuscript
\usepackage{microtype}
\usepackage[hidelinks]{hyperref}
\pagestyle{empty}
\setlength{\parindent}{0pt}
\setlength{\parskip}{\medskipamount}

\begin{document}

\today

\bigskip

Dr.\ Dennis Coates\\
Editor, \emph{Journal of Sports Economics}\\
University of Maryland, Baltimore County

\bigskip

Dear Dr.\ Coates,

Please consider the enclosed manuscript, ``The Limits of Foresight in Coaching Hires: Pre-Hire Predictability of Involuntary Separation and the Internal-Promotion Hazard in the NFL,'' for publication in the \emph{Journal of Sports Economics} as a Research Article.

Head coach turnover in the NFL is frequent and consequential, and the economics of managerial dismissal has consistently traced firing to performance realized once a manager is in the job. That recurring conclusion, that losing leads to firing, offers little to a team deciding whom to hire, because it is expressed in information the employer does not yet hold.

This manuscript asks a different, decision-relevant question: using only what is observable when the contract is signed, how much of an eventual involuntary separation is already foreseeable? I recast the problem as a firing-aware, competing-risks survival analysis over 371 modern-era head-coaching stints (1970 to 2025), modeling the time from hire to firing entirely from pre-hire information, with voluntary departures treated as a competing event and coaches still active at the data boundary censored.

The central result is one of bounded predictability. A cause-specific Cox model, benchmarked against parametric and gradient-boosted alternatives and fit on Cox-native stability-selected covariates, reaches a Harrell concordance of 0.602 and an integrated Brier score of 0.176; the broad on-field record that hiring searches scrutinize forecasts dismissal no better than chance. The one durable pre-hire signal is structural rather than performance-based: coaches promoted from within face a markedly higher firing hazard than external hires (hazard ratio 1.52), a premium that persists after conditioning on observable measures of franchise distress. I read this through the personnel economics of internal versus external hiring, where the direction of the gap is most consistent with a selection account in which an outside candidate is recruited only when he clears a higher bar. The findings bear on how franchises allocate their evaluation effort and on the broader debate over whether managerial turnover is corrective.

The manuscript speaks to the journal's central interests in the labor markets of professional sport, in hiring, and in managerial turnover. It is original, is not under consideration for publication elsewhere, and has not been previously published.

Thank you for considering this manuscript for publication in the \emph{Journal of Sports Economics}.

\bigskip

Sincerely,\\
Jon Williamson\\
Des Moines, Iowa, USA\\
\href{mailto:jonwilliamson@live.com}{jonwilliamson@live.com} \quad \href{mailto:jon@williamsonconsultinggroup.com}{jon@williamsonconsultinggroup.com}

\end{document}
